% ====================================
% CHƯƠNG II - PHẦN 2: YÊU CẦU PHI CHỨC NĂNG & QUY TẮC NGHIỆP VỤ
% Người phụ trách: Người 3
% ====================================


\section{Quy tắc nghiệp vụ}
\label{sec:quy_tac_nghiep_vu}

% Sử dụng danh sách liệt kê để làm gọn mục lục
\begin{itemize}
    \item \textbf{Quy tắc quản lý Series}
    \begin{itemize}
        \item Mỗi Series chỉ thuộc quyền quản lý của một Mangaka.
        \item Một Mangaka có thể quản lý nhiều Series khác nhau.
        \item Series mới phải được Editorial Board xét duyệt trước khi được đưa vào kế hoạch phát triển và xuất bản.
        \item Mỗi Series phải có trạng thái hoạt động để phản ánh tình trạng hiện tại trong quá trình sản xuất và phát hành.
        \item Một Series có thể bao gồm nhiều Chapter.
    \end{itemize}
    \item \textbf{Quy tắc quản lý Chapter}
    \begin{itemize}
        \item Mỗi Chapter phải thuộc về một Series duy nhất.
        \item Một Series có thể bao gồm nhiều Chapter.
        \item Mỗi Chapter phải có số thứ tự riêng trong cùng một Series.
        \item Chapter chỉ được phép phát hành sau khi hoàn tất quá trình kiểm duyệt nội dung.
        \item Mỗi Chapter có thể chứa nhiều Page Manga.
    \end{itemize}
    \item \textbf{Quy tắc quản lý Task}
    \begin{itemize}
        \item Mỗi Task phải được tạo bởi một Mangaka.
        \item Mỗi Task chỉ được giao cho một Assistant tại một thời điểm.
        \item Một Assistant có thể được giao nhiều Task khác nhau.
        \item Mỗi Task phải có thời hạn hoàn thành và trạng thái xử lý.
        \item Chỉ Mangaka mới có quyền tạo, cập nhật hoặc hủy Task.
        \item Tiến độ thực hiện của Task phải được cập nhật và theo dõi trong hệ thống.
    \end{itemize}
    \item \textbf{Quy tắc quản lý Submission}
    \begin{itemize}
        \item Mỗi Submission phải thuộc về một Task cụ thể.
        \item Chỉ Assistant được giao Task mới có quyền gửi Submission tương ứng.
        \item Một Task có thể phát sinh nhiều Submission trong trường hợp cần chỉnh sửa hoặc bổ sung nội dung.
        \item Submission phải được xem xét và đánh giá trước khi được chấp nhận hoàn thành.
        \item Toàn bộ lịch sử Submission phải được lưu trữ trong hệ thống.
    \end{itemize}
    \item \textbf{Quy tắc quản lý Review}
    \begin{itemize}
        \item Mỗi Review phải được thực hiện trên một Submission cụ thể.
        \item Review được thực hiện bởi Tantou Editor.
        \item Nội dung Review phải bao gồm nhận xét, đánh giá hoặc yêu cầu chỉnh sửa.
        \item Kết quả Review là cơ sở để chấp nhận hoặc yêu cầu chỉnh sửa Submission.
        \item Toàn bộ lịch sử Review phải được lưu trữ phục vụ cho việc theo dõi và đánh giá chất lượng nội dung.
    \end{itemize}
    \item \textbf{Quy tắc xuất bản}
    \begin{itemize}
        \item Series mới phải được Editorial Board xét duyệt trước khi đưa vào kế hoạch phát hành.
        \item Chapter chỉ được phát hành sau khi hoàn tất quá trình kiểm duyệt nội dung.
        \item Editorial Board có quyền quyết định tiếp tục hoặc ngừng phát hành một Series.
        \item Thông tin phát hành phải được cập nhật và lưu trữ trong hệ thống.
        \item Các quyết định liên quan đến xuất bản phải dựa trên kết quả đánh giá và dữ liệu hoạt động của Series.
    \end{itemize}
    \item \textbf{Quy tắc xếp hạng}
    \begin{itemize}
        \item Kết quả xếp hạng được xác định dựa trên dữ liệu bình chọn của độc giả.
        \item Mỗi Series chỉ có một vị trí xếp hạng trong cùng một kỳ đánh giá.
        \item Dữ liệu xếp hạng phải được cập nhật theo từng kỳ phát hành hoặc đánh giá.
        \item Kết quả xếp hạng là một trong những căn cứ để Editorial Board đánh giá hiệu quả hoạt động của Series.
        \item Lịch sử xếp hạng phải được lưu trữ để phục vụ công tác thống kê và báo cáo.
    \end{itemize}
    \item \textbf{Quy tắc phân quyền}
    \begin{itemize}
        \item Mỗi người dùng chỉ được gán một vai trò tại một thời điểm.
        \item Mọi thao tác trên hệ thống phải được xác thực thông qua tài khoản người dùng hợp lệ.
        \item Người dùng chỉ được truy cập các chức năng phù hợp với vai trò được cấp.
        \item Admin có quyền quản lý tài khoản người dùng và phân quyền hệ thống.
        \item Mangaka chỉ được quản lý các Series do mình phụ trách.
        \item Assistant chỉ được thực hiện các Task được giao.
        \item Tantou Editor chỉ được thực hiện các chức năng liên quan đến kiểm duyệt và phản hồi nội dung.
        \item Editorial Board chỉ được thực hiện các chức năng liên quan đến xét duyệt, đánh giá, xếp hạng và quyết định xuất bản.
    \end{itemize}
\end{itemize}


\section{Kết luận chương}
\label{sec:ket_luan_chap2}

Chương này đã trình bày chi tiết kết quả phân tích yêu cầu đối với hệ thống Manga Creation Workflow and Publishing Management System. Thông qua việc tìm hiểu nghiệp vụ thực tế, nhóm đã xác định được các bên liên quan (actor) tham gia vào hệ thống bao gồm Admin, Mangaka, Assistant, Tantou Editor và Editorial Board, đồng thời phác thảo quy trình nghiệp vụ tổng quát từ khâu đề xuất Series mới cho đến khi xuất bản và đánh giá xếp hạng.

Bên cạnh đó, chương cũng đã làm rõ các yêu cầu chức năng cốt lõi mà hệ thống cần đáp ứng để phục vụ công việc của từng nhóm người dùng, từ quản lý tài khoản, Series, Chapter, Task đến kiểm duyệt nội dung và xuất bản. Đồng thời, các yêu cầu phi chức năng về hiệu năng, bảo mật, tính sẵn sàng và các quy tắc nghiệp vụ ràng buộc cũng được định nghĩa một cách chặt chẽ.

Những đặc tả và phân tích chi tiết trong chương này sẽ đóng vai trò là nền tảng cơ sở vững chắc để tiến hành thiết kế kiến trúc tổng thể, thiết kế cơ sở dữ liệu và mô hình hóa các biểu đồ UML trong Chương III tiếp theo.
